\section{Abgrenzung und Ausblick}\label{sec:abgrenzung}

\subsection{QAM und Beweisassistenten}

Ein Beweisassistent wie Lean~4~\cite{lean4} prüft den \emph{Inhalt} mathematischer Aussagen: Ein in Lean mit Mathlib~\cite{mathlib}
formalisierter Satz ist bewiesen, relativ zum Kern des Systems. QAM prüft die \emph{Form} einer Akte: welcher Status welcher Aussage zusteht,
gegeben die Status ihrer Prämissen, welche Eingaben sie trägt und in welchem Frame sie gilt. Die beiden Ansätze beantworten verschiedene
Fragen und schließen einander nicht aus.
\begin{center}
\small
\begin{tabular}{@{}p{0.27\textwidth}p{0.32\textwidth}p{0.32\textwidth}@{}}
\toprule
 & Beweisassistent & QAM\\
\midrule
prüft & Gültigkeit eines Beweises & Konsistenz von Status, Eingaben, Frames\\
Eingangsaufwand & vollständige Formalisierung & Deklaration je Aussage\\
offene Aussagen & \texttt{sorry} oder Axiom & erstklassiger Status mit Vererbung\\
physikalische Lesarten & außerhalb des Systems & als \st{reading} registrierbar\\
Rechnungen & im System oder als Orakel & als Auftrag mit Beleg\\
\bottomrule
\end{tabular}
\end{center}
QAM benutzt Lean an zwei Stellen: für seinen eigenen Prüfer und für die Formalisierung der Objektebene und der Maschine. Die Inhalte der
Akten formalisiert es nicht. Eine naheliegende Erweiterung ist, einen Lean-Beweis als Beleg einer Aussage mit Status \st{proved} zu
verknüpfen, so dass dieser Status nicht mehr nur deklariert, sondern belegt ist. Eine solche Verknüpfung besteht derzeit nicht.

\subsection{Verwandte Ansätze}

Die für die Maschine bewiesenen Eigenschaften haben die bekannte Form syntaktischer Typkorrektheit, Erhaltung und Fortschritt~\cite{wright,pierce},
angepasst an eine Maschine, deren Schritte abgewiesen werden können. Die Regel, dass ein Zertifikat seinen Beweis enthält, entspricht der Idee
beweistragenden Codes~\cite{necula}. Das Protokoll der Maschine ist ein Herkunftsnachweis im Sinn des W3C-PROV-Modells~\cite{prov},
reduziert auf Anweisungen und Ausgänge; QAM fügt die Typisierung der aufgezeichneten Objekte hinzu. Vorab-Registrierung ist in Teilen der
empirischen Wissenschaften etabliert~\cite{nosek}; das Projekt überträgt sie auf rechnerische und theoretische Arbeit. Systeme wie
FunSearch~\cite{funsearch} suchen mathematische Konstruktionen mit Sprachmodellen unter automatischer Bewertung, und es gibt Vorschläge für
automatisierte Forschungsabläufe~\cite{aiscientist}. Die hier behandelte Lage ist eine andere: Ein Mensch lenkt ein langes, offenes Projekt,
und das Problem ist nicht die Suche, sondern eine über Monate ehrliche Akte.
\todo{systematische Literaturrecherche zu Provenienzsystemen, Workflow-Sprachen und LLM-gestützter Formalisierung}

\subsection{Weiterentwicklung}

Offen sind: typisierte Frame-Abbildungen, deren Rechtfertigung selbst einen Status trägt, so dass der Hinweis
\texttt{frame-bridge-required} maschinell aufgelöst werden kann; Fortschrittssätze für die speicherverändernden Anweisungen sowie Kopieren
und Zusammenführen von Speichern; ein Nachweis, dass der Python-Adapter die in Lean bewiesene Semantik umsetzt; ein Prüfer, der logische
Prämissenkanten im Forschungstext selbst erkennt, wofür das Inventar der Satzumgebungen die erste Schicht ist; und eine
kollisionsresistente Hashfunktion anstelle von MD5.

\section{Grenzen und Schluss}\label{sec:grenzen}

QAM ersetzt weder einen Beweisassistenten noch die empirische Prüfung. Es prüft nur Deklariertes: Eine als bewiesen deklarierte Aussage
ohne Prämissen wird angenommen, ein unter neuem Namen deklarierter Fitwert wird nicht als Eingabe erkannt, eine falsche Prämisse bleibt
falsch. Frames sind Namen, Abbildungen haben keinen Status, Formeln sind Text. Der Parser des Lean-Prüfers ist nicht verifiziert, die
laufende Maschine ebenfalls nicht. Die Fallstudie ist ein Einzelfall mit einem Autor.

Innerhalb dieser Grenzen leistet QAM eine bestimmte Sache zuverlässig: Es macht sichtbar und überprüfbar, \emph{worauf} jede Aussage ruht, und
verhindert, dass eine offene Bedingung auf dem Weg durch eine lange Kette von Schlüssen verloren geht. Zuses Trennung von Plan und
Ausführung liefert dafür die Architektur: Das Sprachmodell rechnet und formuliert, die Akte legt fest, was gilt. Turings Grenze liefert die
Begründung für den dritten Ausgang: Was nicht entschieden ist, wird nicht geglättet, sondern als Guard geführt, gespeichert und vererbt.
Die Form einer Akte zu sichern macht ihren Inhalt nicht wahr. Aber sie sorgt dafür, dass man sehen kann, welcher Teil des Inhalts noch nicht
gesichert ist.
